% Exemplo: modo matemático no texto (inline) e em destaque, equações numeradas e \eqref
% Compilar: latexmk -pdf modos.tex                 Testado com: TeX Live 2024 ou mais recente
% Cada trecho fica em uma página. Evite $$ ... $$ (sintaxe do TeX puro, com espaçamento errado
% no LaTeX): use \[ ... \] ou os ambientes do amsmath.
\documentclass[a4paper]{article}
\usepackage[T1]{fontenc}
\usepackage{lmodern}
\usepackage[brazilian]{babel}
\usepackage{amsmath}

\begin{document}

%%BEGIN:matinline
No texto, a fórmula fica na linha: $\lim_{x \to \infty} e^{-x} = 0$,
\(x^{1/2} = \sqrt{x}\) e $\int \frac{1}{x}\,dx = \ln\lvert x\rvert + c$. Para
não alterar o espaçamento entre as linhas, o \LaTeX\ reduz frações, limites e
integrais.
%%END:matinline

\clearpage

%%BEGIN:matdestaque
Em destaque, a fórmula fica centralizada em uma linha própria:
\[
  \int \frac{1}{x}\,dx = \ln\lvert x\rvert + c
\]
e, quando precisa ser citada, é numerada com o ambiente \texttt{equation}:
\begin{equation}\label{eq:normal}
  f(x) = \frac{1}{\sigma\sqrt{2\pi}}\, e^{-\frac{(x-\mu)^2}{2\sigma^2}}
\end{equation}
A Equação~\ref{eq:normal} é a densidade da distribuição normal. Com
\verb|\eqref| o número sai entre parênteses: veja \eqref{eq:normal}.
%%END:matdestaque

\clearpage

%%BEGIN:matalign
O ambiente \texttt{align} alinha as linhas no \verb|&| e numera cada uma:
\begin{align}
  f(x) &= (x+a)(x+b) \\
       &= x^2 + (a+b)x + ab
\end{align}
Na versão \texttt{align*} nenhuma linha é numerada; \verb|\notag| retira o
número de uma linha e \verb|\tag{...}| troca o número por um rótulo:
\begin{align}
  f(x) &= (x+a)(x+b) \notag \\
       &= x^2 + (a+b)x + ab \tag{resultado}
\end{align}
%%END:matalign

\end{document}
